\section{Conclusión}
% Que puedo decir con base en lo que obtuve
Este ejercicio en especial refleja la gran herramienta que presentan los AG aplicados en el mundo real, dado que nos enfrentamos a un caso
muy cercano a nosotros siendo este un problema el cual la Universidad en la que estudiamos se enfrenta cada Semestre. Y como es habitual
en estos problemas del mundo real la solucion no es tan sencilla como un problema propuesto para clase, dado que debemos tomar en cuenta muchos
factores, variables y restricciones.

En este caso tomamos solamente una parte de todas las carreras que oferta la UAQ y solo con esto logramos tener una ventana del gran problema que es 
coordinar esto. Sin embargo como se menciona con mayor tiempo y conocimiento de todas las carreras asi como espacios disponibles podemos llegar a una 
respuesta más adecuada. Tambien es importante mencionar el hecho de que en algunos casos se pueden compartir recursos como grupos mezclados, o buscar
un salon libre en otro edificio, asi como algunos casos donde las materias se impartan en linea lo cual libera salones.

Finalmente aterrizando el problema los resultados obtenidos no son perfectos debido a que vemos cruces en algunos puntos, sin embargo esto puede
solucionarse con una poblacion mas variada o incluso prohibir por completo estos cruces en codigo, sin embargo en ambos casos esto hace mas complejo
el problema y a su vez esto se traduce en mayor costo computacional.